\chapter{Measurement and Troubleshooting}
\label{ch:measurement}
This chapter is where the book's models meet a bench, and it introduces almost no new
theory. Every instrument here is a circuit built from parts already analyzed: a
voltmeter is a resistance you add to a Th\'evenin source (\cref{ch:resistive}), a
scope probe is a divider with a pole and a zero (\cref{ch:rc,ch:splane}), a probe's
ground lead is a series RLC resonator (\cref{ch:rlc}), and an antenna analyzer traces
the Smith locus of \cref{ch:lines}. The reason measurement deserves its own chapter is
that these added circuits are usually \emph{invisible on the schematic}, so their
effects get blamed on the circuit under test.

\section{Why Measurement Is a Circuit Problem}
Meters, probes, oscilloscopes, antenna analyzers, and SWR bridges are themselves
circuits. Connecting one changes the circuit unless its impedance is high, low, or
matched in the right way for the measurement---and ``high'' or ``low'' is never
absolute, only relative to the impedance of the node you touch. That is the single
idea this chapter applies four times.

\section{Voltmeter and Ammeter Loading}
A voltmeter with input resistance \(R_{\mathrm{in}}\) placed across a node is simply
a load on that node's Th\'evenin equivalent (\cref{fig:thevenin} in
\cref{ch:resistive}). The measured voltage is therefore what the resulting divider
delivers:
\[
\frac{V_{\text{meas}}}{V_{\text{true}}}
=\frac{R_{\mathrm{in}}}{R_{\mathrm{in}}+R_{Th}} .
\]
Nothing about this is specific to meters---it is \cref{ch:resistive}'s voltage
divider with the instrument playing the role of the lower resistor. The reading is
always \emph{low}, and the error is negligible only when
\(R_{\mathrm{in}}\gg R_{Th}\).

\begin{workedbox}[: how much does the meter disturb?]
Take the divider of \cref{ch:resistive}, whose Th\'evenin resistance was
\(R_{Th}\approx\SI{688}{\ohm}\). A \SI{10}{\mega\ohm} digital multimeter reads
\[
\frac{10\times10^6}{10\times10^6+688}=99.99\%
\]
of the true value---invisible. Even an old analog meter rated
\(\SI{20}{\kilo\ohm\per\volt}\), which on its \SI{10}{\volt} range is only
\SI{200}{\kilo\ohm}, reads \(99.7\%\). Both are fine here.

Now move to a genuinely high-impedance node---a bias network with
\(R_{Th}=\SI{1}{\mega\ohm}\), the sort of thing found at the gate of a field-effect
transistor (FET) or a tube grid:
\[
\text{\SI{10}{\mega\ohm} DMM:}\quad \frac{10}{10+1}=90.9\%
\;\;(\SI{9}{\percent}\text{ low}),
\qquad
\text{\SI{200}{\kilo\ohm} analog:}\quad \frac{0.2}{0.2+1}=16.7\%
\;\;(\SI{83}{\percent}\text{ low}).
\]
The analog meter does not measure that node at all; it replaces it. And notice the
digital meter is \emph{also} wrong by \SI{9}{\percent}---enough to condemn a working
circuit.
\end{workedbox}

An ammeter is the exact dual. It goes \emph{in series}, so its shunt resistance
\(R_{\text{shunt}}\) adds to the loop and reduces the very current it reports:
\[
\frac{I_{\text{meas}}}{I_{\text{true}}}
=\frac{R_{\text{loop}}}{R_{\text{loop}}+R_{\text{shunt}}} ,
\]
which again reads low, and is negligible only when
\(R_{\text{shunt}}\ll R_{\text{loop}}\). A \SI{5}{\ohm} shunt inserted into a
\SI{50}{\ohm} loop costs \SI{9}{\percent} of the current.

\begin{mistakebox}
\textbf{An averaging meter is not an RMS meter.} A cheap AC meter rectifies and
averages, then multiplies by the \(1.111\) form factor of a \emph{sine wave}
(\cref{sec:parseval}). On anything else the calibration is simply wrong---a square
wave reads \SI{11}{\percent} high. A \textbf{true-RMS} meter performs the
root-mean-square integral itself, which by Parseval's relation is the only way to get
a distorted or harmonic-rich waveform right.
\end{mistakebox}

\begin{controlsbox}
\textbf{One rule, stated twice.} A voltmeter must be a much \emph{worse} path than
the node it measures; an ammeter must be a much \emph{better} path than the loop it
measures. Both are the maximum-power-transfer bookkeeping of \cref{ch:resistive} run
in reverse: matching maximizes the power an instrument \emph{takes}, which is
precisely what a measurement must avoid.
\end{controlsbox}

At RF the same logic applies with impedance in place of resistance, and the numbers
get worse: a few centimetres of lead is an inductance, and a probe's tip capacitance
loads a tuned circuit enough to shift its resonance. Touching a high-\(Q\) tank with
a \SI{10}{\pico\farad} probe detunes it---the measurement moves the pole you were
trying to observe.

\section{The Oscilloscope Probe: Pole-Zero Cancellation on the Bench}
\label{sec:probe}
A \(10\!:\!1\) scope probe is the neatest application of \cref{ch:splane} in the whole
book, so it is worth deriving rather than asserting.

The probe contributes a series resistor \(R_1\) with an adjustable capacitor \(C_1\)
across it; the scope input contributes \(R_2\) with its own capacitance \(C_2\)
(including the cable). The result (\cref{fig:probe-schematic}) is a divider whose
two halves are each a parallel RC:
\[
Z_1=\frac{R_1}{1+sR_1C_1},
\qquad
Z_2=\frac{R_2}{1+sR_2C_2}.
\]

\begin{figure}[htbp]
\centering
\begin{circuitikz}[scale=1]
% probe arm: R1 with C1 across it
\draw (0,0) node[left]{\(v_{\mathrm{in}}\)} to[short,o-] (0.6,0);
\draw (0.6,0) to[R,l=\(R_1\)] (3.4,0) coordinate(mid);
\draw (0.6,0) -- (0.6,1.3) to[C,l=\(C_1\)] (3.4,1.3) -- (mid);
% scope input: R2 || C2 to ground
\draw (mid) to[short,-o] (4.6,0) node[right]{\(v_o\)};
\draw (mid) -- (3.4,-0.5);
\draw (2.6,-0.5) to[R,l_=\(R_2\)] (2.6,-2.3);
\draw (4.4,-0.5) to[C,l=\(C_2\)] (4.4,-2.3);
\draw (2.6,-0.5) -- (4.4,-0.5);
\draw (2.6,-2.3) -- (4.4,-2.3);
\node[ground] at (3.5,-2.3){};
\node[GuideBlue] at (0.35,0.42) {\small in};
\node[GuideBlue] at (4.75,0.42) {\small out};
\node[align=right,font=\small] at (1.05,-1.4) {scope\\input};
\end{circuitikz}
\caption{A compensated \(10\!:\!1\) oscilloscope probe. The input is the probe tip
(left) and the output is what the scope actually digitizes (right). The probe supplies
\(R_1\) and the adjustable \(C_1\); the scope and cable supply \(R_2\) and \(C_2\).}
\label{fig:probe-schematic}
\end{figure}

Forming the divider \(H=Z_2/(Z_1+Z_2)\) and clearing denominators,
\[
H(s)=\frac{R_2(1+sR_1C_1)}{R_1(1+sR_2C_2)+R_2(1+sR_1C_1)}
=\underbrace{\frac{R_2}{R_1+R_2}}_{\text{DC divider}}\cdot
\frac{1+sR_1C_1}{1+s\,R_{\parallel}(C_1+C_2)},
\qquad
R_{\parallel}=\frac{R_1R_2}{R_1+R_2}.
\]
Read off the structure: a DC gain of \(R_2/(R_1+R_2)\), one \textbf{zero} at
\(s=-1/(R_1C_1)\), and one \textbf{pole} at \(s=-1/[R_{\parallel}(C_1+C_2)]\). A
frequency-independent divider requires the zero and pole to \emph{cancel}, which
happens when their time constants match:
\[
R_1C_1=R_{\parallel}(C_1+C_2)
\qquad\Longleftrightarrow\qquad
\boxed{\;R_1C_1=R_2C_2\;}
\]
(the middle step is one line of algebra). When that holds, the numerator and
denominator are identical and \(H(s)=R_2/(R_1+R_2)\) exactly---a pure \(10\!:\!1\)
attenuator at every frequency. \emph{Adjusting the trimmer on a probe is placing a
zero on top of a pole.}

\begin{workedbox}[: trimming a probe]
A typical probe has \(R_1=\SI{9}{\mega\ohm}\) into a scope with
\(R_2=\SI{1}{\mega\ohm}\), giving the \(1/10\) DC ratio. If the scope plus cable
present \(C_2=\SI{15}{\pico\farad}\), compensation needs
\[
C_1=\frac{R_2C_2}{R_1}=\frac{15}{9}\approx\SI{1.7}{\pico\farad},
\]
and both time constants then equal \(\SI{15}{\micro\second}\). That is what the
little screw sets.
\end{workedbox}

The classic square-wave test is a direct reading of where the zero sits relative to
the pole (\cref{fig:probe-comp}). For a step input the response of
\(K(1+\tau_zs)/(1+\tau_ps)\) starts at \(K\tau_z/\tau_p\) and relaxes to \(K\):
\begin{description}
\item[\(\tau_z<\tau_p\) (\(C_1\) too small)] the trace starts low and creeps up---the
familiar \emph{rounded} leading edge. The zero is left of the pole, so it cannot
cancel the pole's lag.
\item[\(\tau_z=\tau_p\)] flat. Zero on pole, no dynamics at all.
\item[\(\tau_z>\tau_p\) (\(C_1\) too large)] the trace overshoots and decays---a
\emph{peaked} edge, because the zero now dominates at high frequency.
\end{description}

\begin{figure}[htbp]
\centering
\includegraphics[width=\textwidth]{probe_compensation}
\caption{Probe compensation as pole-zero cancellation, for
\(R_1=\SI{9}{\mega\ohm}\), \(R_2=\SI{1}{\mega\ohm}\), \(C_2=\SI{15}{\pico\farad}\).
\emph{Left:} step response for three trimmer settings---rounded, flat, and peaked.
\emph{Right:} the same three cases as pole-zero plots (the three rows are offset
vertically for clarity; all the poles and zeros are real and lie on the \(\sigma\)
axis). Compensation is achieved when the zero lands exactly on the pole, which is the
condition \(R_1C_1=R_2C_2\).}
\label{fig:probe-comp}
\end{figure}

\begin{exambox}
The probe is worth remembering as the concrete answer to ``what does a zero
\emph{do}?'' A pole alone rolls the response off; a zero placed on top of it cancels
that roll-off exactly---the extreme case of the slope bookkeeping in
\cref{sec:cascadeadd}, where a pole bends the magnitude down by
\SI{20}{\decibel}/decade and a zero bends it back up by the same amount. The same
cancellation, deliberately mistuned, is how oscilloscope engineers build peaking
networks---and the same idea, applied to a transfer function instead of a divider,
is what \cref{ch:filters} calls placing a notch.
\end{exambox}

\section{Ground-Lead Inductance: an Accidental Resonator}
The other half of probe technique is the ground lead, and it is a pure application of
\cref{ch:rlc}. A wire has roughly \SI{20}{\nano\henry} per inch of inductance. Clip a
six-inch ground lead onto a circuit and you have added \(L\approx\SI{120}{\nano\henry}\)
in series with the probe's own tip capacitance, \(C\approx\SI{10}{\pico\farad}\)---a
series LC loop with almost no resistance to damp it. Its resonance is
\cref{ch:rlc}'s formula, unchanged:
\[
f_0=\frac{1}{2\pi\sqrt{LC}}
=\frac{1}{2\pi\sqrt{(120\times10^{-9})(10\times10^{-12})}}
\approx\SI{145}{\mega\hertz}.
\]

\begin{physicalbox}
\textbf{The overshoot is yours, not the circuit's.} Because that loop has very little
resistance, its \(Q\) is high---the pole pair sits close to the \(\jj\omega\) axis---so
any fast edge shock-excites it and it \emph{rings} at about \SI{145}{\mega\hertz}. A
scope shows overshoot and ringing on a signal that may have had none. The cure follows
straight from \(f_0=1/(2\pi\sqrt{LC})\): shorten the ground path. A one-inch spring
tip cuts \(L\) to \SI{20}{\nano\henry} and pushes the resonance to about
\SI{356}{\mega\hertz}, out of the band of interest. You are not damping the resonance,
you are moving its poles up in frequency until they no longer matter.
\end{physicalbox}

\section{Ground Loops: KVL Around the Return Path}
\label{sec:groundloops}
The ground-lead resonance above is one failure of an assumption the schematic makes
everywhere: that ``ground'' is a single node at zero volts. It is not. A ground
conductor is a conductor, with resistance and---more importantly at RF---inductance,
so current flowing in it develops a voltage, and two points labelled ``ground'' are
at the same potential only if no current flows between them.

Make that concrete with KVL. Suppose two pieces of equipment are joined by a signal
cable \emph{and} both are grounded to the mains. Then there is a closed loop:
out along the cable shield, through one chassis, down its safety ground, along the
building wiring, and back up the other. Any current circulating in that loop---
induced by a nearby transformer at \SI{60}{\hertz}, or by a transmitter at
RF---develops a voltage along the shield, and that voltage adds directly to the
signal, because the receiving end measures its input \emph{relative to} its own
ground. This is a \textbf{ground loop}, and \SI{60}{\hertz} hum is its signature.

The cure follows from the loop equation rather than from folklore. Since the problem
is a circulating current in a closed path, either break the path or shrink its
impedance:
\begin{description}
\item[Single-point (star) grounding] bond every chassis to \emph{one} common point
with its own short conductor, so no two ground paths form a loop at all. This is the
structural fix.
\item[Bonding] where a loop is unavoidable, make its enclosed impedance as low as
possible with short, wide straps---a strap has less inductance than a round wire of
the same length, which matters because it is \(\omega L\) and not \(R\) that
dominates at RF.
\item[Common-mode choking] add series impedance to the loop current specifically,
with a ferrite over the cable (\cref{sec:commonmode}). This attacks the current
rather than the path.
\end{description}

\begin{physicalbox}
\textbf{Why long ground leads fail twice.} A ground conductor a significant
fraction of a wavelength long is not a short at all: by \cref{ch:lines} its input
impedance is \(\jj Z_0\tan\beta\ell\), which passes through \emph{infinity} at a
quarter wavelength. So a ``ground'' wire can present a high impedance and develop
large RF voltages at its open end---which is how an earthed chassis comes to bite
the operator, and why RF grounding is about electrical length rather than about
digging a deeper rod. It is the same tangent that made a stub a resonator in
\cref{sec:stubs}, arriving here as a hazard instead of a component.
\end{physicalbox}

\section{SWR Bridge and Antenna Analyzer}
An SWR bridge separates the forward and reflected waves of \cref{ch:lines} and reports
their ratio as SWR. An antenna analyzer goes further and sweeps frequency, reporting
impedance and SWR across a band---which means its display \emph{is} the Smith locus of
\cref{fig:smith-antenna}. Everything the analyzer tells you was derived in
\cref{ch:lines,ch:rlc}: the frequency where the trace crosses the real axis is
\(\omega_0=1/\sqrt{LC}\) for the antenna's series-RLC model, the distance from center
is \(|\Gamma|\), and the width of the low-SWR region is \(BW\approx f_0/Q\).

\begin{mistakebox}
\textbf{SWR measured at the shack is not SWR at the antenna.} A \emph{lossless} line
rotates \(\Gamma\) on a constant-\(|\Gamma|\) circle (\cref{ch:lines}), so it changes
the impedance you see but not the SWR. A \emph{lossy} line attenuates the reflected
wave on its way back, shrinking \(|\Gamma|\) and making the match look better than it
is---quantitatively, \(|\Gamma_{\mathrm{in}}|=|\Gamma_L|\ee^{-2\alpha\ell}\)
(\cref{sec:lineloss}). Through \SI{3}{\decibel} of matched-line loss the reflected
wave makes a \SI{6}{\decibel} round trip, halving \(|\Gamma|\):
\[
\text{antenna SWR } 5\!:\!1 \;\;(|\Gamma|=0.67)
\quad\longrightarrow\quad
|\Gamma_{\mathrm{in}}|=0.33,
\quad
\text{shack SWR } \approx 2\!:\!1 .
\]
A comfortable-looking \(2\!:\!1\) at the transmitter can therefore hide a badly
mismatched antenna---and the power is not saved, it is dissipated in the coax. Measure
at the antenna feed point when it matters.
\end{mistakebox}

\section{Decibels in Measurement}
The decibel definitions of \cref{sec:decibels} apply unchanged, but on a bench they
usually carry a \emph{reference}, which converts a ratio into an absolute level:
\[
P_{\dB}=10\log_{10}\frac{P}{P_{\mathrm{ref}}},
\qquad
V_{\dB}=20\log_{10}\frac{V}{V_{\mathrm{ref}}}.
\]
Always check which reference is meant. \textbf{dBW} references \(\SI{1}{\watt}\) and
\textbf{dBm} references \(\SI{1}{\milli\watt}\), so the two differ by exactly
\SI{30}{\decibel} (\(\SI{1}{\watt}=\SI{30}{dBm}\)). A bare ``dB'' with no reference
is a ratio---gain, loss, or return loss---and not a level at all. The most important
referenced level in receiver work is the thermal noise floor of
\cref{sec:noisefloor}, \SI{-174}{dBm} per hertz.

\subsection{Adding It All Up: Link Budgets and S Units}
\label{sec:sunits}
The reason referenced decibels are worth the bookkeeping is that they turn a chain of
multiplications into a sum. \Cref{sec:cascadeadd} established this for cascaded
transfer functions; the same arithmetic applied to a whole radio path is a
\textbf{link budget}. Start from the transmitter power as a level, then add every gain
and subtract every loss along the way:
\[
P_{\text{rx}}\,[\mathrm{dBm}]
=P_{\text{tx}}-L_{\text{feed}}+G_{\text{tx}}
-L_{\text{path}}+G_{\text{rx}}-L_{\text{feed,rx}} .
\]
The \textbf{link margin} is then simply how much is left over above the receiver's
minimum usable level---a subtraction, not a ratio, because both are already levels.

\begin{workedbox}[: a link budget in one line]
A \SI{100}{\watt} transmitter (\SI{50}{dBm}) feeds \SI{100}{\foot} of RG-213 at
\SI{144}{\mega\hertz}, so \SI{2}{\decibel} of line loss (\cref{sec:lineloss}), into
an antenna with \SI{10}{\decibel} of gain. The path costs \SI{130}{\decibel} and the
receiving station has a \SI{3}{\decibel} antenna and \SI{1}{\decibel} of feed loss:
\[
50-2+10-130+3-1=\SI{-70}{dBm}.
\]
If the receiver needs \SI{-110}{dBm} to be usable, the link margin is
\SI{40}{\decibel}---comfortable. Note what the arithmetic makes obvious: the
\SI{2}{\decibel} of feed line and the \SI{10}{\decibel} of antenna gain are directly
comparable, and trading a better cable against a better antenna is a matter of
counting decibels.
\end{workedbox}

Receiver reports use their own referenced scale. One \textbf{S unit} is defined as
\SI{6}{\decibel}, which is a factor of four in power (\(10^{6/10}\approx3.98\)). So a
signal reported one S unit stronger needs four times the transmitter power, all else
equal, and ``\SI{20}{\decibel} over S9'' means a hundred times the power of an S9
signal. The awkwardness of the scale is exactly that it is logarithmic while the
power behind it is not: going from S8 to S9 is the same \SI{6}{\decibel}---the same
\(4\times\)---as going from S3 to S4.

\begin{mistakebox}
Do not read S units as a linear signal-strength scale. Each unit is
\SI{6}{\decibel}, so two S units is \SI{12}{\decibel} (\(16\times\) the power), not
twice one unit's worth. And above S9 the scale is abandoned entirely in favor of
``\(n\) dB over S9,'' precisely because the S numbers would otherwise run out.
\end{mistakebox}

\section{Troubleshooting Pattern}
The chapters of this book supply a checklist, because every one of them rests on
assumptions that a real circuit can violate:
\begin{enumerate}
\item Identify the expected DC operating point.
\item Check whether the signal path is AC-coupled, DC-coupled, filtered, or matched.
\item Estimate the node's impedance \emph{before} attaching an instrument, and confirm
your meter or probe is far enough off it (\cref{sec:probe}).
\item Compare the measured gain, attenuation, phase, or SWR against the simple model.
\item If they differ, ask which assumption failed, and note that each one belongs to a
specific chapter: the \emph{ideal element} and \emph{lumped} assumptions of
\cref{ch:foundations}; the \emph{sinusoidal steady state} of \cref{ch:ac}, which says
nothing about transients; the \emph{no-loading} assumption this chapter is about; and
the \emph{linearity} that all of Part~II requires and that a saturating amplifier or a
mixer (\cref{ch:active}) deliberately breaks.
\end{enumerate}

\begin{exambox}
Measurement questions reward recognizing what the instrument actually loads or
reports: voltage, current, power, impedance, SWR, frequency, or phase. When a
measurement disagrees with theory, the instrument is a circuit element too---and it is
usually the one that was left off the schematic.
\end{exambox}
